\title{Leveraging Large Language Models for Automated Dialogue Analysis}

\begin{abstract}
Developing high-performing dialogue systems benefits from the automatic identification of undesirable behaviors in system responses. However, detecting such behaviors remains challenging, as it draws on a breadth of general knowledge and understanding of conversational practices. Although recent research has focused on building specialized classifiers for detecting specific dialogue behaviors, the behavior coverage is still incomplete and there is a lack of testing on real-world human-bot interactions. This paper investigates the ability of a state-of-the-art large language model (LLM), \linebreak ChatGPT-3.5, to perform dialogue behavior detection for nine categories in real human-bot dialogues. We aim to assess whether ChatGPT can match specialized models and approximate human performance, thereby reducing the cost of behavior detection tasks. Our findings reveal that neither specialized models nor ChatGPT have yet achieved satisfactory results for this task, falling short of human performance. Nevertheless, ChatGPT shows promising potential and often outperforms specialized detection models. We conclude with an in-depth examination of the prevalent shortcomings of ChatGPT, offering guidance for future research to enhance LLM capabilities.
\end{abstract}

\section{Introduction}

One crucial aspect of developing high-performing dialogue systems is the automated identification of errors in system responses. These errors can result from various behaviors, including incorrect information retrieval or illogical semantics (Figure~\ref{fig:example_dialogue}).
Identifying such errors enhances dialogue system development and complements dialogue-level evaluation methods by providing finer-grained metrics for comparison \cite{finch:23}.

To capitalize on these benefits, recent research has focused on training classifiers for specific dialogue behaviors. While certain behaviors have received considerable attention, this is not the case for all pertinent dialogue behaviors. Furthermore, most datasets for training are produced by annotating human-human dialogues \cite{sharma:20}, perturbing human responses \cite{gupta:22}, or crafting post-hoc responses \cite{nie:21}. As a result, such datasets may not reflect human-bot interactions, rendering them less suitable for classifier development.

Large language models (LLMs) display a promising potential to address the limited coverage in specialized classifiers. LLMs have demonstrated competitive performance across various natural language processing (NLP) tasks without finetuning \cite{kocon:23}. 

Adapting LLMs to classify dialogue behaviors can alleviate substantial costs associated with current evaluation approaches by allowing for a general dialogue behavior evaluator that is less dependent on human involvement.

Although there is much effort towards open-sourcing competitive LLMs, OpenAI's ChatGPT remains the most successful LLM to date \cite{wang:23}. Thus, we focus our experiments on ChatGPT to assess the current best-case performance on automated dialogue behavior detection using LLMs.

With its wide accessibility and low costs, ChatGPT provides a practical and straightforward platform for automating dialogue behavior detection, if its proves successful.

 To this end, our work focuses on two main objectives:

\begin{enumerate}

\item To determine whether or not ChatGPT can match the performance of state-of-the-art specialized behavior classifiers.
\item To assess the extent to which ChatGPT can approximate human-level performance in behavior classification using real human-bot dialogues.
\end{enumerate}

 Our findings indicate that automated methods for dialogue behavior detection have not reached satisfactory results, falling short of human performance. However, ChatGPT showcases compelling results comparative to or often better than specialized models. To facilitate further advancements, we conduct an in-depth analysis to identify the prevalent errors and shortcomings of ChatGPT. This analysis provides valuable insights, highlighting key areas to be targeted to enhance the performance of LLMs in dialogue behavior detection for future work.   We release our code and data at \url{https://github.com/emorynlp/GPT-ABCEval}.

\section{Related Work}
\label{sec:rw}

ChatGPT has shown promising performance on many NLP tasks, especially for text classification \cite{gilardi:23, kocon:23, zhu:23}. In addition, GPT models, including ChatGPT and InstructGPT, have been used to produce high-quality dyadic dialogues \cite{kim:22, zhan:23} and have been shown to correlate highly with human annotators when evaluating the overall quality of empathetic dialogues \cite{svikhnushina:23}. However, ChatGPT still exhibits limitations as \citet{chan:23} show that ChatGPT struggles with fine-grained dialogue understanding, reporting poor performance on classifying discourse structure and utterance relations. 

To the best of our knowledge, no prior research has explored the use of any GPT model as a behavior classifier for chatbot responses. Instead, previous work has focused on the development of specialized dialogue behavior classifiers, as discussed in this section.

\subsection{Contradiction Detection}
\label{sec:rw_speaker_contradiction}

Although much work focuses on dialogue contradictions in the context of a given bot persona \cite{zhang:18, welleck:19, kim:20, song:20, shuster:22}, there has been some work on a more general sense of contradictions, including NLI models targeting self-context contradictions \cite{li:21, nie:21}, inconsistency detectors using domain-specific attribute-value classifiers \cite{shi:21}, and context summarization to encourage consistency in response generation \cite{xu_beyond:22, xu_long:22}. Notably, these existing approaches to contradiction detection fail to address partner contradictions.

There is also a lack of work on general commonsense contradiction detection for dialogue responses. To the best of our knowledge, \citet{ghazarian:23} is the only work that focuses explicitly on capturing commonsense qualities of dialogue responses. They propose a method for calculating continuous event commonsense alignment scores for dialogue responses using similarity calculations with the outputs of an event extraction model and generative commonsense model. However, such continuous scores cannot be immediately applied to commonsense contradiction detection without further modifications (e.g. learned thresholding, classification head, etc.).

\subsection{Claim Verification}
\label{sec:rw_claim_verification}

There are a variety of approaches taken for claim verification in dialogue, including question-answering \cite{honovich:21} and trained classifiers \cite{dziri_begin:22}. \citet{dziri_begin:22} find that trained classifiers perform the best, although they still lag behind human performance. Some works focus on claim verification for question-response pairs only \cite{wang:22}, whereas others target multi-turn dialogues, producing annotated datasets including FaithDial \cite{dziri_faithdial:22}, BEGIN \cite{dziri_begin:22}, and DialFact \cite{gupta:22}. Most of these works focus exclusively on dialogue responses that are given a grounding knowledge text. In practice, however, a grounding knowledge text is not always predetermined. \citet{gupta:22} propose a pipeline for claim verification that includes a knowledge retrieval stage rather than assuming it is provided.

\subsection{Empathy}
\label{sec:rw_empathy}

Human judges are commonly used when evaluating the degree of empathy exhibited in a dialogue response \cite{zhong:20,sabour:22,qian:23}. There has also been some work on developing empathetic response and question taxonomies, although these are only applied in small-scale or synthetic settings \cite{welivita:20, svikhnushina:22}. Most applicably, \citet{sharma:20} collect EPITOME, a dataset of 10K interactions from Reddit and Talklife (a mental health forum) that are annotated with the strength of their expression of three empathetic mechanisms: reactions, interpretations, explorations. Some recent dialogue works have used EPITOME-trained classifiers in their approaches \cite{zheng:21, majumder:22} or for automatic evaluation \cite{kim:21, lee:22}.

\subsection{Coherence}
\label{sec:rw_coherence}

Research on detecting incoherent behaviors, such as redundancy and irrelevancy, is limited. Most works perturb dialogue responses to artificially construct incoherence examples \cite{xu:21, zhang:21, ghazarian:22}, which may not produce representative examples. On the other hand, \citet{mehri:20} derive a response's relevancy score from the probabilities of manually designed future indicator utterances but found little correlation with human judgments. In addition, detection of response redundancy is underexplored, despite some works addressing token repetition \cite{li:20, xi:21}. Perhaps most relevant, the Dialogue Breakdown Detection Challenge (DBDC) aims to identify contextually inappropriate bot responses that hinder conversation continuation \cite{higashinaka:19}. Various classifiers have been proposed for this challenge \cite{ng:20,lin:22}, with observations suggesting coherence issues as a dominant cause of breakdowns.

\section{ABC-Eval Dataset}
\label{sec:dataset} 

We use the ABC-Eval dataset from \citet{finch:23} as the behavior detection benchmark. This dataset contains 400 open-domain human-bot dialogues collected between university students and one of four chatbots: BlenderBot2, Blenderbot using DECODE reranking, Emora, and Bart-FiD-RAG. For each bot response in each dialogue, human annotators labeled whether or not a specific dialogue behavior was present. These turn-level binary annotations were collected using crowdworking annotators on the SurgeHQ platform,\footnote{\url{https://www.surgehq.ai}} who were trained on three curated conversations to accurately identify each dialogue behavior before being accepted into the annotation project. For example, in Figure \ref{fig:example_dialogue}, the three bot responses are labeled \texttt{1}, \texttt{0}, \texttt{0} for the behavior \texttt{incorrect fact (!Fac)} and are labeled \texttt{0}, \texttt{0}, \texttt{1} for the behavior  \texttt{self contradiction (!Sel)}.

In this work, we take 1,634 bot responses from 108 dialogues that received two rounds of human annotations, and focus on the nine dialogue behaviors that \citet{finch:23} found as the most informative for capturing dialogue quality (Table~\ref{tab:behavior_labels}).

\section{Specialized Behavior Detection Models}
\label{sec:models}

In this section, we present state-of-the-art models designed to classify labels that closely align with six of the dialogue behaviors in Table~\ref{tab:behavior_labels}: \texttt{Emp}, \texttt{!Emp}, \texttt{!Fac}, \texttt{!Sel}, \texttt{Ign}, and \texttt{!Rel}.
Note that no existing models are available for predicting\linebreak \texttt{!Com}, \texttt{!Par}, and \texttt{Red} so there are no viable comparisons to our LLM approach for them (Section~\ref{sec:gpt_llm}).

\paragraph{FaithCritic (FC)} 

Following \citet{gupta:22}, we build a claim verification pipeline for a dialogue response $r$. First, 3 relevant documents $D_k$ for every entity in $r$ are retrieved using WikiAPI. Then, a BERT model trained on the Wizard of Wikipedia (WoW) knowledge-response pairs \cite{dinan:19} selects the top-10 evidence sentences $S_e$ from $D_k$. To distinguish whether a response makes a factual claim or not, the lexical overlap between $r$ and $S_e$ is estimated, optimized on the ABC-Eval training conversations. Finally, a RoBERTa model trained on Faith-Critic, a dataset of human-annotated faithful and unfaithful evidence-response pairs derived from the WoW \cite{dziri_faithdial:22}, is applied to those responses that make factual claims. As a result, responses that are predicted unfaithful to any evidence $e \in S_e$ are labeled as \texttt{!Fac}.

\paragraph{S2T2}  

S2T2 is a semi-supervised student-teacher training framework using two teachers, one trained on the gold data and the other trained on perturbed gold data under a [MASK] replacement, to incorporate self-supervised data augmentation into the model training \cite{lin:22}. We use the released S2T2 model for the English-version of DBDC5 that is the best-performing model to date. We use S2T2 as identifying \texttt{Ign} and \texttt{!Rel} labels, since it is not trained to distinguish between them.

\paragraph{DECODE (DEC)} 

We use the released RoBERTa classification model trained on DECODE to label \texttt{!Sel}. DECODE contains human-written contradictory and non-contradictory dialogue responses with respect to the current speaker's previous utterances in the dialogue \cite{nie:21}.

\paragraph{EPITOME (EPI)} 

A RoBERTa-based bi-encoder classification model for each empathetic communication mechanism is trained from the publicly available Reddit portion of the EPITOME dataset \cite{sharma:20}. Predictions of weak or strong expressions of any of the three mechanisms are considered as \texttt{Emp}. Predictions of no expression for all mechanisms are considered as \texttt{!Emp}.

\section{LLM-based Behavior Detection}
\label{sec:gpt_llm}

For LLM-based dialogue behavior detection, we use OpenAI’s \textit{gpt-turbo-3.5-301} (henceforth, \linebreak ChatGPT). Similar to the specialized models (Section~\ref{sec:models}), ChatGPT is tasked with classifying a single behavior at a time.

Following the human annotator training process for ABC-Eval, we use the three training conversations for each label as our prompt engineering testbed.

This section highlights key decisions of our prompt engineering process.

\paragraph{Instruction Finetuning}

During prompt engineering, it became apparent that the instructions designed for human annotators (Section~\ref{sec:dataset}) were not suitable as ChatGPT instructions. We iteratively refined the instructions such that ChatGPT's mistakes on the training conversations were reduced. This involved removing instructions ChatGPT appeared to misunderstand as well as adding additional behavior details and specifications.

\paragraph{Utterance Focus} 

We discovered that when ChatGPT was instructed to label each bot turn given the entire dialogue, the resulting classifications often focused on only a subset of the bot responses. To ensure consistent and robust labeling for every bot utterance, our final prompt provides the dialogue history paired with the next bot response as the target utterance to be labeled.

\paragraph{In-context Examples}

We also tried including the examples provided to human annotators by \citet{finch:23} as in-context examples in the prompts. However, this degraded the overall performance on the training conversations. It appears that the examples optimized for improving \textit{human} annotations do not translate well to ChatGPT's performance.

\paragraph{Creativity}

We conducted experiments involving several \textit{temperature} parameters and observed high instability in the classifications for the same inputs when the temperature was increased. Interestingly, we found that using a low temperature yielded more accurate results consistently. Thus, we opted to use a \textit{temperature} of 0 for maximum reproducibility in our classifications. Similar findings have been reported by \citet{gilardi:23} and \citet{reiss:23}.

\paragraph{Final Prompt}

Table~\ref{tab:final_prompt} illustrates an example of the final prompt, in particular for the \texttt{Red} behavior\footnote{Due to spacing constraints, Table \ref{tab:final_prompt} contains minor discrepancies with the actual prompts in this work (Appendix \ref{app:full_prompt_example}).}. 
Each behavior is accompanied by its own eliciting question and description, which respectively fill in the \textit{Behavior Question} (\textbf{Q}) and \textit{Behavior Definition} (\textbf{D}) containers of the prompt. For labeling a particular context-response example, the historical turns from the context and the next target response fill in the \textit{Dialogue Context} (\textbf{C}) and \textit{Target Utterance} (\textbf{U}) containers. The final label for a behavior is parsed from the line produced by ChatGPT that begins with the header `\textit{Decision:}', where the value `yes' indicates a positive label and any other value indicates a negative label. Appendix \ref{app:prompt_values} provides the questions and definitions used for all behaviors. 

\section{Evaluation}
\label{sec:evaluation_metrics}

To evaluate the detection capability of the models in Sections~\ref{sec:models} and \ref{sec:gpt_llm}, we compare their performance against that of human annotators. For this, we take the set of doubly annotated conversations in ABC-Eval as our evaluation set (108 dialogues), and apply each model to the bot responses (1,634 utterances) to obtain the predicted labels.

\subsection{Metrics}

To assess the degree to which automated methods can approximate human judgment for a particular dialogue behavior, we measure the accuracy of the binary labels predicted by automated methods with respect to the binary labels provided by the human annotators. In addition, we calculate both the F1-score for the positive occurrences of each dialogue behavior and for the negative occurrences of each dialogue behavior, in order to obtain a more fine-grained picture of the performance.

Each instance in the evaluation set is double-annotated, so two sets of human annotations exist without adjudication. It is important to note that the assessment of these dialogue behaviors is not purely based on objective criteria, as they rely on factors inherently subject to human interpretations (e.g., commonsense contradiction, irrelevance). With this in mind, to better capture the aggregate nature of identifying dialogue behaviors, the final score for each metric is measured by averaging results across the double human annotations, where $e$ is the metric (either accuracy or F1-score), $o_m$ is the model outputs, and $o_{h1}$ and $o_{h2}$ are the human labels from annotation round 1 and 2, respectively: $$ e_{final} = \frac{1}{2}(e(o_m, o_{h1}) + e(o_m, o_{h2})) $$ To assess human performance, we measure the F1 score and accuracy by comparing the two human annotation sets. Finally, the statistical significance between outputs of models and humans, and between outputs of the specialized models and ChatGPT, is estimated using McNemar's Test with significance level of $0.05$. Testing is performed by treating each human annotation set as ground-truth.\footnote{The other human annotation set relative to the one being treated as ground-truth is used as human output.}

\subsection{Results \& Discussion}
\label{sec:results}

 Table~\ref{tab:results} indicates the ongoing challenge of dialogue behavior detection for automated models. Across all labels, human judges are significantly more stable than the models. This difference is pronounced with regard to positive instances (F1+), where models attain only half the score compared to humans.

Interestingly, ChatGPT exhibits comparable performance with several specialized classifiers. In the case of \texttt{!Fac}, ChatGPT outperforms FaithCritic (FC) in every aspect and achieves performance closer to humans. For \texttt{!Emp} and \texttt{!Rel}, ChatGPT shows similar performance on F1- and accuracy, and even better performance on F1+, as their classifiers. Considering that ChatGPT is not finetuned for these tasks, these results are highly encouraging.

Although ChatGPT is seemingly outperformed by S2T2 on \texttt{Ign}, this is primarily due to the prediction of negative cases. When analyzing the positive cases, ChatGPT gives much higher recall yet similar precision compared to S2T2\footnote{Precision and recall provided in Appendix \ref{app:full_results}.}.
In practice, positive case detection is more impactful, implying that ChatGPT has an advantage in real-world applications.

Furthermore, although ChatGPT faces significant challenges in detecting positive cases of \texttt{Emp}, EPITOME (EPI) does not perform much better. Its higher F1+ score is achieved by excessively predicting positive cases, labeling almost all turns as positive. This overprediction impairs its overall performance, allowing ChatGPT to outperform it when considering all cases as reflected in accuracy.

The only behavior for which ChatGPT appears to be beaten by the specialized classifier is against DECODE (DEC) for \texttt{!Sel}. However, the difference in performance is only slight overall. 

Notably, ChatGPT shows promising accuracy and negative F1 (F1-) to humans for the three behaviors for which specialized models are not available: \texttt{!Com}, \texttt{!Par}, and \texttt{Red}.
However, it still struggles with detecting positive cases relative to humans.

\section{ChatGPT Error Analysis}
\label{sec:gpt_error_analysis}

We perform an error analysis of ChatGPT's predictions of dialogue behaviors to better understand its limitations. For each dialogue behavior, we select 40 instances where ChatGPT and humans disagree, and examine the reasoning provided by ChatGPT prior to its final decision (\textbf{\texttt{[R]}}; see examples below).
Table~\ref{tab:error_analysis} presents a set of dialogue characteristics and ChatGPT predispositions that highlight common mistakes made by ChatGPT across multiple dialogue behaviors.

\subsection{Context Management}

The predominant cause of ChatGPT's errors is its inability to focus on the pertinent parts of the dialogue and response. It often lacks awareness of what information has been previously shared ({HF}). In Figure~\ref{fig:history-forgetfulness}, the earlier response that already answers ``\textit{Are you older than your brother}'' is missed.

 Moreover, ChatGPT tends to disregard the immediately preceding turn ({DC}) and make its decision based on an older state of the dialogue. In Figure~\ref{fig:disassociated-context}, it overlooks the preceding question about favorite colors. Similarly, ChatGPT often misses ideas shared within a response, instead latching on the very last part of the dialogue ({SA}). In Figure~\ref{fig:selective-attention}, ChatGPT fails to notice the emotional mimicry conveyed by S2 in response to S1's opinion about Nicki Minaj.

 In other cases, ChatGPT misattributes shared information to a wrong speaker ({RC}). In Figure~\ref{fig:role-confusion}, it recalls that S2 indicated a return to normalcy, when in fact, it was S1 who made that claim.

\subsection{Instruction Following}

In addition, a frequent issue is that ChatGPT strays from the provided behavior definitions. In Figure \ref{fig:definition-mismatch}, esoteric knowledge is improperly considered as commonsense. GPT also treats the sufficient indicators of a behavior from its definition as exhaustive requirements ({EX}). In Figure~\ref{fig:exhaustive}, ChatGPT criticizes S2 for not indicating an emotion and offering support, despite S2's validation of S1's experience.

\subsection{World Model}

Another major issue is ChatGPT's poor understanding of common world events and human experiences.

ChatGPT frequently reveals a limited understanding of the relationship among concepts within a dialogue. For example, it overlooks elaborations on previous points, considering them too semantically similar (SR). In Figure~\ref{fig:semantic-relatedness}, the opinion about the cat's characteristics is actually unique information relative to the context.

 ChatGPT also often criticizes typical conversational practices ({CN}). In Figure~\ref{fig:conversation-norms}, it considers topical introductions at the start of a dialogue as irrelevant.

 In addition, ChatGPT faces challenges in comprehending the plausibility of co-occurring events, beliefs, and experiences ({ME}). In Figure~\ref{fig:mutual-exclusion}, it depicts a lack of understanding that it is implausible to have an opinion about the music of an artist if one has no prior experience with that artist's music.

ChatGPT also demonstrates a large degree of general inexperience with common phenomenon or situations in the world, which can lead to it harboring untrue facts about the world or misunderstanding nuances of situations ({IN}). In Figure~\ref{fig:inexperience}, ChatGPT's lack of commonsense is highlighted by its acceptance of well-wishes for a non-existent trip.

\section{Recommendations}
\label{sec:recommendations}

Given the compelling performance for many dialogue behaviors observed in this work, ChatGPT is a promising direction for behavior classification.

For one, it is worth noting that ChatGPT boasts extreme cost-efficiency relative to humans. Where ChatGPT costs \$0.02 on average to provide labels for a single behavior for one dialogue in this work, the average cost for human annotation ranges from \$0.29 to \$1.96 depending on the behavior (Table~\ref{tab:cost_comparison} in Appendix~\ref{sec:appendix_cost}). Since even specialized classifiers rely on human annotations for training creation, they also end up being quite costly to maintain. 

Furthermore, the results of our error analysis reveal a large degree of systematicity behind ChatGPT's reasoning mistakes across many of the behaviors. Correcting these common mistakes is likely to further improve its performance to a noticeable degree. We next discuss mitigation strategies of these identified issues to aid in future work.

\paragraph{Context Management} Providing the complete dialogue history may hinder ChatGPT's ability to attend to the salient content due to information overload. To address this, we highlight two strategies:

\begin{itemize}
    \item \textit{Windowed Context}: instead of providing the entire history, truncate the context to \textit{k} previous turns. This would directly restrict the decision-making to the immediate context, which is important for behaviors that depend on accurate recency identification, including \texttt{!Rel}, \texttt{Ign}, \texttt{!Emp}, and \texttt{Emp}.
    \item \textit{Turn Pairing}: perform the labeling relative to each historical turn segment independently, rather than a contiguous context. This would enable explicit and focused comparisons to smaller segments of the history that could aid behaviors that require such precision, including \texttt{!Sel}, \texttt{!Par}, and \texttt{Red}.
\end{itemize}

\paragraph{In-Context Learning Examples} Given the identified mistake types, it becomes more straightforward to compose useful in-context learning examples that are tailored to optimizing ChatGPT. Examples of those mistake types that are related to ChatGPT misunderstanding the nuances of a behavior (e.g. MD, SR, CN, ME, EX) could be taken from a held-out set of conversations, which would prime ChatGPT to avoid such reasoning.

\section{Limitations}

Although ChatGPT is a high-performing, widely accessible, and affordable LLM at the time of writing, there are considerations towards the long-term applicability of the results found in this work due to the ChatGPT infrastructure. Since ChatGPT is not open-source and is only accessible through a paid API, there is less detailed understanding of its training and model design. In addition, this access method for ChatGPT also results in less user control over potential model changes and even model deprecation over time. As such, further studies could assess the applicability of other language models to the task of dialogue behavior detection to mitigate these concerns, and we leave this to future work.

Furthermore, it should be noted that the errors made by ChatGPT may not necessarily align with those made by alternative open-source language models, or even future versions of ChatGPT itself. However, it may still be useful to be mindful of the prominent problems encountered with ChatGPT while using other LLMs. These identified phenomena play a crucial role in language comprehension and reasoning overall and could also present challenges for other models, although the extent of their impact remains to be explored.

\section{Conclusion}
\label{sec:conclusion}

Although automated methods for dialogue behavior classification remain a challenging task, this work finds that ChatGPT-3.5 presents promising potential to reduce the gap between model and human performance. ChatGPT's ability to provide competitive behavior classification against specialized classifiers without necessitating finetuning or human annotation across a variety of dialogue behaviors gives rise to a low-cost, multi-task evaluator model. The systematicity behind the common mistakes observed for ChatGPT reveal concrete steps for future improvements that will improve behavior classification performance, including strategies for context management and better understanding of situational nuances. We look forward to future advancements in behavior classification that leverage ChatGPT's unique capabilities.

\end{document}
